%%%%%%%%%%%%%%%%%%%%%%%%%%%%%%%%%%%%%%%%%%%%%%%%%%%%%%%%%%%%
% 2.1 ELEMENTARY NUMBER THEORY
% The Composition of Advanced Mathematics
%%%%%%%%%%%%%%%%%%%%%%%%%%%%%%%%%%%%%%%%%%%%%%%%%%%%%%%%%%%%

\section{Elementary Number Theory}
\label{sec:elementary-number-theory}

\prerequisite{
Arithmetic foundations, number systems, elementary set theory, functions,
logic, and methods of proof from Chapter 0. Familiarity with elementary
algebra from \cref{sec:elementary-algebra} is also assumed.
}

\learningobjectives{
By the end of this section, the reader should be able to:
\begin{itemize}
    \item explain the central questions studied in elementary number theory;
    \item work with divisibility and divisibility notation;
    \item prove basic properties of divisibility;
    \item apply the Division Algorithm;
    \item compute greatest common divisors and least common multiples;
    \item use the Euclidean Algorithm and Extended Euclidean Algorithm;
    \item express greatest common divisors as linear combinations;
    \item identify relatively prime integers;
    \item work with prime and composite numbers;
    \item apply the Fundamental Theorem of Arithmetic;
    \item use prime factorizations to compute gcds and lcms;
    \item work with congruences and modular arithmetic;
    \item solve linear congruences;
    \item compute multiplicative inverses modulo an integer;
    \item solve systems using the Chinese Remainder Theorem;
    \item apply Fermat's Little Theorem and Euler's Theorem;
    \item use the Euler phi function;
    \item recognize elementary Diophantine equations;
    \item solve linear Diophantine equations;
    \item understand quadratic residues at an introductory level;
    \item recognize connections between number theory, algebra,
    computation, and cryptography.
\end{itemize}
}

%%%%%%%%%%%%%%%%%%%%%%%%%%%%%%%%%%%%%%%%%%%%%%%%%%%%%%%%%%%%
\subsection{What Is Number Theory?}
\label{subsec:what-is-number-theory}
%%%%%%%%%%%%%%%%%%%%%%%%%%%%%%%%%%%%%%%%%%%%%%%%%%%%%%%%%%%%

Number theory is the study of the arithmetic structure and properties of
integers.

The integers

\[ \Z=\{\ldots,-3,-2,-1,0,1,2,3,\ldots\} \]

appear elementary, but their internal structure leads to some of the deepest
questions in mathematics.

Typical questions include:

\begin{itemize}
    \item When does one integer divide another?
    \item How can two integers share common factors?
    \item How are integers built from prime numbers?
    \item When do integer equations have integer solutions?
    \item What happens when integers are considered modulo another integer?
    \item How can arithmetic properties be used in computation and cryptography?
\end{itemize}

\begin{conceptbox}
Elementary number theory studies integers through divisibility, prime
factorization, greatest common divisors, congruences, and integer-valued
equations.

Although the objects are elementary, the methods include algebra,
algorithms, induction, contradiction, and other forms of mathematical proof.
\end{conceptbox}

%%%%%%%%%%%%%%%%%%%%%%%%%%%%%%%%%%%%%%%%%%%%%%%%%%%%%%%%%%%%
\subsection{Divisibility}
\label{subsec:number-theory-divisibility}
%%%%%%%%%%%%%%%%%%%%%%%%%%%%%%%%%%%%%%%%%%%%%%%%%%%%%%%%%%%%

Divisibility is one of the fundamental relations in number theory.

\begin{definitionbox}{Divisibility}
Let \(a,b\in\Z\). We say that \(a\) \emph{divides} \(b\) if there exists an
integer \(k\) such that

\[ b=ak. \]

This is written

\[ a\mid b. \]
\end{definitionbox}

The symbol

\[ \mid \]

is called the \emph{divides symbol}.

The statement

\[ a\mid b \]

is read

\begin{center}
``\(a\) divides \(b\).''
\end{center}

For example,

\[ 4\mid20 \]

because

\[ 20=4(5). \]

Similarly,

\[ -3\mid15 \]

because

\[ 15=(-3)(-5). \]

However,

\[ 4\nmid18. \]

The symbol

\[ \nmid \]

is read ``does not divide.''

%%%%%%%%%%%%%%%%%%%%%%%%%%%%%%%%%%%%%%%%%%%%%%%%%%%%%%%%%%%%
\subsection{Divisor and Multiple}
\label{subsec:divisor-multiple}
%%%%%%%%%%%%%%%%%%%%%%%%%%%%%%%%%%%%%%%%%%%%%%%%%%%%%%%%%%%%

If

\[ a\mid b, \]

then \(a\) is called a \emph{divisor} or \emph{factor} of \(b\), while \(b\)
is called a \emph{multiple} of \(a\).

For example, since

\[ 6\mid42, \]

the integer \(6\) is a divisor of \(42\), and \(42\) is a multiple of \(6\).

%%%%%%%%%%%%%%%%%%%%%%%%%%%%%%%%%%%%%%%%%%%%%%%%%%%%%%%%%%%%
\subsection{Basic Properties of Divisibility}
\label{subsec:divisibility-properties}
%%%%%%%%%%%%%%%%%%%%%%%%%%%%%%%%%%%%%%%%%%%%%%%%%%%%%%%%%%%%

Divisibility behaves predictably under multiplication and addition.

\begin{theorem}[Basic Divisibility Properties]
Let \(a,b,c,m,n\in\Z\).

\begin{enumerate}
    \item \(a\mid a\).
    \item \(1\mid a\).
    \item If \(a\mid b\) and \(b\mid c\), then \(a\mid c\).
    \item If \(a\mid b\), then \(a\mid bc\).
    \item If \(a\mid b\) and \(a\mid c\), then
    \[ a\mid(b+c) \qquad\text{and}\qquad a\mid(b-c). \]
    \item If \(a\mid b\) and \(a\mid c\), then
    \[ a\mid(mb+nc). \]
\end{enumerate}
\end{theorem}

The final property is especially important.

The expression

\[ mb+nc \]

is called an \emph{integer linear combination} of \(b\) and \(c\).

%%%%%%%%%%%%%%%%%%%%%%%%%%%%%%%%%%%%%%%%%%%%%%%%%%%%%%%%%%%%
\subsection{Proof of a Divisibility Property}
\label{subsec:divisibility-property-proof}
%%%%%%%%%%%%%%%%%%%%%%%%%%%%%%%%%%%%%%%%%%%%%%%%%%%%%%%%%%%%

Suppose

\[ a\mid b \qquad\text{and}\qquad a\mid c. \]

By definition, there exist integers \(r,s\) such that

\[ b=ar,\qquad c=as. \]

Let \(m,n\in\Z\). Then

\[ mb+nc=m(ar)+n(as)=a(mr+ns). \]

Since

\[ mr+ns\in\Z, \]

the definition of divisibility gives

\[ \boxed{a\mid(mb+nc).} \]

This illustrates an important number-theoretic proof strategy: begin with the
definition of divisibility and rewrite the relevant integers as multiples.

%%%%%%%%%%%%%%%%%%%%%%%%%%%%%%%%%%%%%%%%%%%%%%%%%%%%%%%%%%%%
\subsection{The Division Algorithm}
\label{subsec:division-algorithm}
%%%%%%%%%%%%%%%%%%%%%%%%%%%%%%%%%%%%%%%%%%%%%%%%%%%%%%%%%%%%

Ordinary integer division is formalized by the Division Algorithm.

\begin{theorem}[Division Algorithm]
Let \(a\in\Z\) and let \(d\in\Z\) with \(d>0\). Then there exist unique
integers \(q\) and \(r\) such that

\[ a=dq+r,\qquad 0\leq r<d. \]
\end{theorem}

Here:

\begin{itemize}
    \item \(a\) is the \emph{dividend};
    \item \(d\) is the \emph{divisor};
    \item \(q\) is the \emph{quotient};
    \item \(r\) is the \emph{remainder}.
\end{itemize}

\begin{centereddiagram}
\begin{tikzpicture}[
    box/.style={draw,rounded corners,minimum width=2.5cm,minimum height=1cm}
]
\node[box] (a) at (0,0) {Dividend \(a\)};
\node[box] (d) at (4,0) {Divisor \(d\)};
\node[box] (q) at (0,-2) {Quotient \(q\)};
\node[box] (r) at (4,-2) {Remainder \(r\)};

\node at (2,-3.2) {\(\displaystyle a=dq+r,\qquad 0\leq r<d\)};

\draw[-{Latex[length=2.5mm]}] (a)--(q);
\draw[-{Latex[length=2.5mm]}] (d)--(r);
\end{tikzpicture}
\end{centereddiagram}

%%%%%%%%%%%%%%%%%%%%%%%%%%%%%%%%%%%%%%%%%%%%%%%%%%%%%%%%%%%%
\subsection{Worked Example: Division Algorithm}
\label{subsec:worked-division-algorithm}
%%%%%%%%%%%%%%%%%%%%%%%%%%%%%%%%%%%%%%%%%%%%%%%%%%%%%%%%%%%%

\begin{workedexamplebox}{Finding Quotient and Remainder}

Express \(137\) in the form

\[ 137=12q+r,\qquad 0\leq r<12. \]

\begin{tabularx}{\textwidth}{@{}>{\raggedright\arraybackslash}p{0.42\textwidth}X@{}}
Find the largest multiple of \(12\) not exceeding \(137\). &
\(12(11)=132\)\\
Find the remainder. &
\(137-132=5\)\\
Write the Division Algorithm form. &
\(137=12(11)+5\)
\end{tabularx}

\begin{answerbox}
\[ \boxed{q=11,\qquad r=5.} \]
\end{answerbox}
\end{workedexamplebox}

%%%%%%%%%%%%%%%%%%%%%%%%%%%%%%%%%%%%%%%%%%%%%%%%%%%%%%%%%%%%
\subsection{Division with Negative Integers}
\label{subsec:division-negative-integers}
%%%%%%%%%%%%%%%%%%%%%%%%%%%%%%%%%%%%%%%%%%%%%%%%%%%%%%%%%%%%

The remainder condition

\[ 0\leq r<d \]

must still be satisfied when the dividend is negative.

For example, dividing \(-17\) by \(5\) gives

\[ -17=5(-4)+3. \]

Therefore

\[ q=-4,\qquad r=3. \]

Writing

\[ -17=5(-3)-2 \]

is algebraically correct, but it is not the Division Algorithm form because

\[ r=-2<0. \]

%%%%%%%%%%%%%%%%%%%%%%%%%%%%%%%%%%%%%%%%%%%%%%%%%%%%%%%%%%%%
\subsection{Greatest Common Divisor}
\label{subsec:greatest-common-divisor}
%%%%%%%%%%%%%%%%%%%%%%%%%%%%%%%%%%%%%%%%%%%%%%%%%%%%%%%%%%%%

Two integers may have several common divisors.

For example, the positive divisors of \(24\) are

\[ 1,2,3,4,6,8,12,24, \]

while the positive divisors of \(36\) are

\[ 1,2,3,4,6,9,12,18,36. \]

Their positive common divisors are

\[ 1,2,3,4,6,12. \]

The largest is \(12\).

\begin{definitionbox}{Greatest Common Divisor}
Let \(a,b\in\Z\), not both zero. The \emph{greatest common divisor} of \(a\)
and \(b\), written

\[ \gcd(a,b), \]

is the greatest positive integer that divides both \(a\) and \(b\).
\end{definitionbox}

Thus

\[ \gcd(24,36)=12. \]

%%%%%%%%%%%%%%%%%%%%%%%%%%%%%%%%%%%%%%%%%%%%%%%%%%%%%%%%%%%%
\subsection{Relatively Prime Integers}
\label{subsec:relatively-prime}
%%%%%%%%%%%%%%%%%%%%%%%%%%%%%%%%%%%%%%%%%%%%%%%%%%%%%%%%%%%%

\begin{definitionbox}{Relatively Prime}
Two integers \(a\) and \(b\) are \emph{relatively prime}, or \emph{coprime},
if

\[ \gcd(a,b)=1. \]
\end{definitionbox}

For example,

\[ \gcd(8,15)=1, \]

so \(8\) and \(15\) are relatively prime.

Notice that neither integer needs to be prime.

%%%%%%%%%%%%%%%%%%%%%%%%%%%%%%%%%%%%%%%%%%%%%%%%%%%%%%%%%%%%
\subsection{Least Common Multiple}
\label{subsec:least-common-multiple}
%%%%%%%%%%%%%%%%%%%%%%%%%%%%%%%%%%%%%%%%%%%%%%%%%%%%%%%%%%%%

\begin{definitionbox}{Least Common Multiple}
For nonzero integers \(a\) and \(b\), the \emph{least common multiple},
written

\[ \operatorname{lcm}(a,b), \]

is the smallest positive integer divisible by both \(a\) and \(b\).
\end{definitionbox}

For example,

\[ \operatorname{lcm}(12,18)=36. \]

For positive integers \(a\) and \(b\),

\[ \boxed{\gcd(a,b)\operatorname{lcm}(a,b)=ab.} \]

More generally,

\[ \boxed{\operatorname{lcm}(a,b)=\frac{|ab|}{\gcd(a,b)}.} \]

%%%%%%%%%%%%%%%%%%%%%%%%%%%%%%%%%%%%%%%%%%%%%%%%%%%%%%%%%%%%
\subsection{The Euclidean Algorithm}
\label{subsec:euclidean-algorithm}
%%%%%%%%%%%%%%%%%%%%%%%%%%%%%%%%%%%%%%%%%%%%%%%%%%%%%%%%%%%%

Listing all divisors becomes inefficient for large integers.

The Euclidean Algorithm computes greatest common divisors using repeated
division.

The key observation is

\[ \gcd(a,b)=\gcd(b,r) \]

whenever

\[ a=bq+r. \]

Repeatedly applying the Division Algorithm eventually produces remainder
zero.

The last nonzero remainder is the gcd.

\begin{centereddiagram}
\begin{tikzpicture}[
    node distance=8mm,
    box/.style={draw,rounded corners,align=center,minimum width=4.2cm,
    minimum height=0.8cm}
]
\node[box] (one) {\(a=bq_1+r_1\)};
\node[box,below=of one] (two) {\(b=r_1q_2+r_2\)};
\node[box,below=of two] (three) {\(r_1=r_2q_3+r_3\)};
\node[below=5mm of three] (dots) {\(\vdots\)};
\node[box,below=5mm of dots] (last) {\(r_{n-2}=r_{n-1}q_n+r_n\)};
\node[box,below=of last] (zero) {\(r_{n-1}=r_nq_{n+1}+0\)};

\draw[-{Latex[length=2.5mm]}] (one)--(two);
\draw[-{Latex[length=2.5mm]}] (two)--(three);
\draw[-{Latex[length=2.5mm]}] (three)--(dots);
\draw[-{Latex[length=2.5mm]}] (dots)--(last);
\draw[-{Latex[length=2.5mm]}] (last)--(zero);

\node[right=8mm of zero] {\(\gcd(a,b)=r_n\)};
\end{tikzpicture}
\end{centereddiagram}

%%%%%%%%%%%%%%%%%%%%%%%%%%%%%%%%%%%%%%%%%%%%%%%%%%%%%%%%%%%%
\subsection{Worked Example: Euclidean Algorithm}
\label{subsec:worked-euclidean-algorithm}
%%%%%%%%%%%%%%%%%%%%%%%%%%%%%%%%%%%%%%%%%%%%%%%%%%%%%%%%%%%%

\begin{workedexamplebox}{Finding a Greatest Common Divisor}

Find

\[ \gcd(252,198). \]

\begin{tabularx}{\textwidth}{@{}>{\raggedright\arraybackslash}p{0.42\textwidth}X@{}}
Divide \(252\) by \(198\). &
\(252=198(1)+54\)\\
Divide \(198\) by \(54\). &
\(198=54(3)+36\)\\
Divide \(54\) by \(36\). &
\(54=36(1)+18\)\\
Divide \(36\) by \(18\). &
\(36=18(2)+0\)\\
The last nonzero remainder is the gcd. &
\(\gcd(252,198)=18\)
\end{tabularx}

\begin{answerbox}
\[ \boxed{\gcd(252,198)=18.} \]
\end{answerbox}
\end{workedexamplebox}

%%%%%%%%%%%%%%%%%%%%%%%%%%%%%%%%%%%%%%%%%%%%%%%%%%%%%%%%%%%%
\subsection{Bézout's Identity}
\label{subsec:bezout-identity}
%%%%%%%%%%%%%%%%%%%%%%%%%%%%%%%%%%%%%%%%%%%%%%%%%%%%%%%%%%%%

A gcd can always be expressed as an integer linear combination.

\begin{theorem}[Bézout's Identity]
If \(a,b\in\Z\), not both zero, then there exist integers \(x\) and \(y\)
such that

\[ ax+by=\gcd(a,b). \]
\end{theorem}

The integers \(x\) and \(y\) are called \emph{Bézout coefficients}.

The name Bézout is commonly pronounced approximately ``bay-ZOO.''

%%%%%%%%%%%%%%%%%%%%%%%%%%%%%%%%%%%%%%%%%%%%%%%%%%%%%%%%%%%%
\subsection{The Extended Euclidean Algorithm}
\label{subsec:extended-euclidean-algorithm}
%%%%%%%%%%%%%%%%%%%%%%%%%%%%%%%%%%%%%%%%%%%%%%%%%%%%%%%%%%%%

The Extended Euclidean Algorithm determines Bézout coefficients by working
backward through the Euclidean Algorithm.

Using the previous example,

\[ 252=198+54, \]

\[ 198=3(54)+36, \]

\[ 54=36+18. \]

Work backward:

\[ 18=54-36. \]

Since

\[ 36=198-3(54), \]

we obtain

\[ 18=54-[198-3(54)]=4(54)-198. \]

Since

\[ 54=252-198, \]

we obtain

\[ 18=4(252-198)-198. \]

Therefore

\[ 18=4(252)-5(198). \]

Hence

\[ \boxed{252(4)+198(-5)=18.} \]

%%%%%%%%%%%%%%%%%%%%%%%%%%%%%%%%%%%%%%%%%%%%%%%%%%%%%%%%%%%%
\subsection{A Characterization of Coprime Integers}
\label{subsec:coprime-bezout}
%%%%%%%%%%%%%%%%%%%%%%%%%%%%%%%%%%%%%%%%%%%%%%%%%%%%%%%%%%%%

Bézout's Identity immediately gives an important criterion.

\begin{theorem}
Integers \(a\) and \(b\) are relatively prime if and only if there exist
integers \(x,y\) such that

\[ ax+by=1. \]
\end{theorem}

This theorem will later become essential when determining whether an integer
has a multiplicative inverse modulo another integer.

%%%%%%%%%%%%%%%%%%%%%%%%%%%%%%%%%%%%%%%%%%%%%%%%%%%%%%%%%%%%
\subsection{Euclid's Lemma}
\label{subsec:euclids-lemma}
%%%%%%%%%%%%%%%%%%%%%%%%%%%%%%%%%%%%%%%%%%%%%%%%%%%%%%%%%%%%

\begin{theorem}[Euclid's Lemma]
If \(a,b,c\in\Z\), with

\[ a\mid bc \]

and

\[ \gcd(a,b)=1, \]

then

\[ a\mid c. \]
\end{theorem}

\begin{proof}
Since

\[ \gcd(a,b)=1, \]

Bézout's Identity gives integers \(x,y\) satisfying

\[ ax+by=1. \]

Multiply by \(c\):

\[ acx+bcy=c. \]

Clearly,

\[ a\mid acx. \]

Since \(a\mid bc\),

\[ a\mid bcy. \]

Therefore \(a\) divides their sum, so

\[ a\mid c. \]
\end{proof}

%%%%%%%%%%%%%%%%%%%%%%%%%%%%%%%%%%%%%%%%%%%%%%%%%%%%%%%%%%%%
\subsection{Prime Numbers}
\label{subsec:prime-numbers}
%%%%%%%%%%%%%%%%%%%%%%%%%%%%%%%%%%%%%%%%%%%%%%%%%%%%%%%%%%%%

Prime numbers are the multiplicative building blocks of the positive
integers.

\begin{definitionbox}{Prime Number}
An integer \(p>1\) is \emph{prime} if its only positive divisors are

\[ 1\qquad\text{and}\qquad p. \]
\end{definitionbox}

Examples include

\[ 2,3,5,7,11,13,17,19,23,\ldots \]

The number \(2\) is the only even prime number.

%%%%%%%%%%%%%%%%%%%%%%%%%%%%%%%%%%%%%%%%%%%%%%%%%%%%%%%%%%%%
\subsection{Composite Numbers}
\label{subsec:composite-numbers}
%%%%%%%%%%%%%%%%%%%%%%%%%%%%%%%%%%%%%%%%%%%%%%%%%%%%%%%%%%%%

\begin{definitionbox}{Composite Number}
An integer \(n>1\) is \emph{composite} if it is not prime.
\end{definitionbox}

Equivalently, \(n\) is composite if

\[ n=ab \]

for integers satisfying

\[ 1<a<n,\qquad 1<b<n. \]

For example,

\[ 12=3(4), \]

so \(12\) is composite.

\begin{warningbox}
The integer \(1\) is neither prime nor composite.

This convention is necessary for the uniqueness statement in the
Fundamental Theorem of Arithmetic.
\end{warningbox}

%%%%%%%%%%%%%%%%%%%%%%%%%%%%%%%%%%%%%%%%%%%%%%%%%%%%%%%%%%%%
\subsection{Prime Divisibility}
\label{subsec:prime-divisibility}
%%%%%%%%%%%%%%%%%%%%%%%%%%%%%%%%%%%%%%%%%%%%%%%%%%%%%%%%%%%%

Euclid's Lemma takes an especially important form for primes.

\begin{theorem}
If \(p\) is prime and

\[ p\mid ab, \]

then

\[ p\mid a\qquad\text{or}\qquad p\mid b. \]
\end{theorem}

This property distinguishes prime numbers from general composite integers.

For example,

\[ 6\mid2(3), \]

but

\[ 6\nmid2,\qquad6\nmid3. \]

Thus the corresponding statement fails for composite divisors.

%%%%%%%%%%%%%%%%%%%%%%%%%%%%%%%%%%%%%%%%%%%%%%%%%%%%%%%%%%%%
\subsection{The Fundamental Theorem of Arithmetic}
\label{subsec:fundamental-theorem-arithmetic}
%%%%%%%%%%%%%%%%%%%%%%%%%%%%%%%%%%%%%%%%%%%%%%%%%%%%%%%%%%%%

\begin{theorem}[Fundamental Theorem of Arithmetic]
Every integer \(n>1\) can be written as a product of prime numbers.

Moreover, this factorization is unique except for the order of the prime
factors.
\end{theorem}

For example,

\[ 360=2^3\cdot3^2\cdot5. \]

Although we could reorder the factors, the primes and their exponents are
uniquely determined.

\begin{conceptbox}
Prime numbers play the role of multiplicative building blocks:

\[
\boxed{
\text{positive integers}
\longleftrightarrow
\text{products of primes}
}
\]
\end{conceptbox}

%%%%%%%%%%%%%%%%%%%%%%%%%%%%%%%%%%%%%%%%%%%%%%%%%%%%%%%%%%%%
\subsection{Canonical Prime Factorization}
\label{subsec:canonical-prime-factorization}
%%%%%%%%%%%%%%%%%%%%%%%%%%%%%%%%%%%%%%%%%%%%%%%%%%%%%%%%%%%%

A positive integer can be written in the form

\[ n=p_1^{\alpha_1}p_2^{\alpha_2}\cdots p_k^{\alpha_k}, \]

where

\[ p_1<p_2<\cdots<p_k \]

are distinct primes and

\[ \alpha_i\in\N. \]

This is called the \emph{canonical prime factorization} of \(n\).

For example,

\[ 756=2^2\cdot3^3\cdot7. \]

%%%%%%%%%%%%%%%%%%%%%%%%%%%%%%%%%%%%%%%%%%%%%%%%%%%%%%%%%%%%
\subsection{GCD and LCM from Prime Factorization}
\label{subsec:gcd-lcm-prime-factorization}
%%%%%%%%%%%%%%%%%%%%%%%%%%%%%%%%%%%%%%%%%%%%%%%%%%%%%%%%%%%%

Suppose

\[ a=\prod_{i=1}^{k}p_i^{\alpha_i},\qquad b=\prod_{i=1}^{k}p_i^{\beta_i}, \]

where zero exponents may be inserted for missing primes.

Then

\[ \boxed{\gcd(a,b)=\prod_{i=1}^{k}p_i^{\min(\alpha_i,\beta_i)}} \]

and

\[ \boxed{\operatorname{lcm}(a,b)=\prod_{i=1}^{k}p_i^{\max(\alpha_i,\beta_i)}.} \]

%%%%%%%%%%%%%%%%%%%%%%%%%%%%%%%%%%%%%%%%%%%%%%%%%%%%%%%%%%%%
\subsection{Worked Example: GCD and LCM by Factorization}
\label{subsec:worked-gcd-lcm-factorization}
%%%%%%%%%%%%%%%%%%%%%%%%%%%%%%%%%%%%%%%%%%%%%%%%%%%%%%%%%%%%

\begin{workedexamplebox}{Using Prime Factorization}

Find the gcd and lcm of

\[ 360\qquad\text{and}\qquad504. \]

\begin{tabularx}{\textwidth}{@{}>{\raggedright\arraybackslash}p{0.42\textwidth}X@{}}
Factor \(360\). &
\(360=2^3\cdot3^2\cdot5\)\\
Factor \(504\). &
\(504=2^3\cdot3^2\cdot7\)\\
Take the smaller exponents for the gcd. &
\(\gcd(360,504)=2^3\cdot3^2=72\)\\
Take the larger exponents for the lcm. &
\(\operatorname{lcm}(360,504)=2^3\cdot3^2\cdot5\cdot7=2520\)
\end{tabularx}

\begin{answerbox}
\[ \boxed{\gcd(360,504)=72,\qquad\operatorname{lcm}(360,504)=2520.} \]
\end{answerbox}
\end{workedexamplebox}

%%%%%%%%%%%%%%%%%%%%%%%%%%%%%%%%%%%%%%%%%%%%%%%%%%%%%%%%%%%%
\subsection{Infinitely Many Primes}
\label{subsec:infinitely-many-primes}
%%%%%%%%%%%%%%%%%%%%%%%%%%%%%%%%%%%%%%%%%%%%%%%%%%%%%%%%%%%%

One of the oldest results in mathematics states that the prime numbers never
end.

\begin{theorem}[Euclid]
There are infinitely many prime numbers.
\end{theorem}

\begin{proof}
Suppose, for contradiction, that there are only finitely many primes:

\[ p_1,p_2,\ldots,p_n. \]

Construct

\[ N=p_1p_2\cdots p_n+1. \]

Since \(N>1\), it has a prime divisor \(q\).

By assumption, \(q\) must be one of

\[ p_1,\ldots,p_n. \]

Therefore

\[ q\mid p_1p_2\cdots p_n. \]

Also,

\[ q\mid N. \]

Hence \(q\) divides the difference

\[ N-p_1p_2\cdots p_n=1. \]

But no prime divides \(1\).

This is a contradiction.

Therefore there must be infinitely many primes.
\end{proof}

%%%%%%%%%%%%%%%%%%%%%%%%%%%%%%%%%%%%%%%%%%%%%%%%%%%%%%%%%%%%
\subsection{Testing for Primality}
\label{subsec:elementary-primality-testing}
%%%%%%%%%%%%%%%%%%%%%%%%%%%%%%%%%%%%%%%%%%%%%%%%%%%%%%%%%%%%

Suppose \(n>1\) is composite.

Then

\[ n=ab \]

for integers \(a,b>1\).

At least one of these factors must satisfy

\[ a\leq\sqrt n \qquad\text{or}\qquad b\leq\sqrt n. \]

Therefore, to test whether \(n\) is prime, it is enough to test divisibility
by primes not exceeding

\[ \sqrt n. \]

For example, to test \(97\), we only need primes up to

\[ \sqrt{97}<10. \]

Thus test

\[ 2,3,5,7. \]

None divides \(97\), so \(97\) is prime.

%%%%%%%%%%%%%%%%%%%%%%%%%%%%%%%%%%%%%%%%%%%%%%%%%%%%%%%%%%%%
\subsection{Congruence}
\label{subsec:congruence}
%%%%%%%%%%%%%%%%%%%%%%%%%%%%%%%%%%%%%%%%%%%%%%%%%%%%%%%%%%%%

Modular arithmetic studies integers according to their remainders.

\begin{definitionbox}{Congruence Modulo \(n\)}
Let \(a,b,n\in\Z\) with \(n>0\). We say that \(a\) is \emph{congruent to}
\(b\) modulo \(n\) if

\[ n\mid(a-b). \]

We write

\[ a\equiv b\pmod n. \]
\end{definitionbox}

The symbol

\[ \equiv \]

is called the \emph{congruence symbol} in this context and is read
``is congruent to.''

The notation

\[ a\equiv b\pmod n \]

is read

\begin{center}
``\(a\) is congruent to \(b\) modulo \(n\).''
\end{center}

%%%%%%%%%%%%%%%%%%%%%%%%%%%%%%%%%%%%%%%%%%%%%%%%%%%%%%%%%%%%
\subsection{Congruence and Remainders}
\label{subsec:congruence-remainders}
%%%%%%%%%%%%%%%%%%%%%%%%%%%%%%%%%%%%%%%%%%%%%%%%%%%%%%%%%%%%

Two integers are congruent modulo \(n\) precisely when they have the same
remainder after division by \(n\).

For example,

\[ 17\equiv5\pmod{12} \]

because

\[ 17-5=12. \]

Likewise,

\[ 38\equiv2\pmod{12} \]

because

\[ 38-2=36=12(3). \]

Thus modulo \(12\),

\[ 2,14,26,38,50,\ldots \]

all represent the same remainder class.

%%%%%%%%%%%%%%%%%%%%%%%%%%%%%%%%%%%%%%%%%%%%%%%%%%%%%%%%%%%%
\subsection{Congruence Classes}
\label{subsec:congruence-classes}
%%%%%%%%%%%%%%%%%%%%%%%%%%%%%%%%%%%%%%%%%%%%%%%%%%%%%%%%%%%%

For \(a\in\Z\), define

\[ [a]_n=\{b\in\Z:b\equiv a\pmod n\}. \]

This is the \emph{congruence class} or \emph{residue class} of \(a\) modulo
\(n\).

For example, modulo \(5\), every integer belongs to exactly one of

\[ [0]_5,\ [1]_5,\ [2]_5,\ [3]_5,\ [4]_5. \]

These classes partition \(\Z\).

%%%%%%%%%%%%%%%%%%%%%%%%%%%%%%%%%%%%%%%%%%%%%%%%%%%%%%%%%%%%
\subsection{The Integers Modulo \(n\)}
\label{subsec:integers-modulo-n}
%%%%%%%%%%%%%%%%%%%%%%%%%%%%%%%%%%%%%%%%%%%%%%%%%%%%%%%%%%%%

The set of congruence classes modulo \(n\) is written

\[ \Z/n\Z. \]

It may be represented as

\[ \Z/n\Z=\{[0]_n,[1]_n,\ldots,[n-1]_n\}. \]

This notation was developed algebraically in Chapter 1 through quotient
structures. Here we focus on its arithmetic interpretation.

%%%%%%%%%%%%%%%%%%%%%%%%%%%%%%%%%%%%%%%%%%%%%%%%%%%%%%%%%%%%
\subsection{A Modular Number Line}
\label{subsec:modular-number-line}
%%%%%%%%%%%%%%%%%%%%%%%%%%%%%%%%%%%%%%%%%%%%%%%%%%%%%%%%%%%%

Modulo \(5\), integer values repeat every five units.

\begin{centereddiagram}
\begin{tikzpicture}[scale=0.9]
\draw[-{Latex[length=2.5mm]}] (-6.5,0)--(6.5,0);

\foreach \x in {-6,-5,...,6}{
    \draw (\x,0.12)--(\x,-0.12);
    \node[below] at (\x,-0.12) {\small \x};
}

\node[above] at (-5,0.15) {\small \(0\)};
\node[above] at (0,0.15) {\small \(0\)};
\node[above] at (5,0.15) {\small \(0\)};

\draw[<->,dashed] (-5,0.8)--(0,0.8);
\node[above] at (-2.5,0.8) {\(5\) units};

\draw[<->,dashed] (0,0.8)--(5,0.8);
\node[above] at (2.5,0.8) {\(5\) units};
\end{tikzpicture}
\end{centereddiagram}

Numbers separated by a multiple of \(5\) are congruent modulo \(5\).

%%%%%%%%%%%%%%%%%%%%%%%%%%%%%%%%%%%%%%%%%%%%%%%%%%%%%%%%%%%%
\subsection{Properties of Congruence}
\label{subsec:congruence-properties}
%%%%%%%%%%%%%%%%%%%%%%%%%%%%%%%%%%%%%%%%%%%%%%%%%%%%%%%%%%%%

Congruence behaves much like equality under arithmetic operations.

If

\[ a\equiv b\pmod n \]

and

\[ c\equiv d\pmod n, \]

then

\[ a+c\equiv b+d\pmod n, \]

\[ a-c\equiv b-d\pmod n, \]

and

\[ ac\equiv bd\pmod n. \]

For every positive integer \(k\),

\[ a^k\equiv b^k\pmod n. \]

These properties allow large arithmetic expressions to be reduced modulo
\(n\) during a calculation.

%%%%%%%%%%%%%%%%%%%%%%%%%%%%%%%%%%%%%%%%%%%%%%%%%%%%%%%%%%%%
\subsection{Worked Example: Modular Arithmetic}
\label{subsec:worked-modular-arithmetic}
%%%%%%%%%%%%%%%%%%%%%%%%%%%%%%%%%%%%%%%%%%%%%%%%%%%%%%%%%%%%

\begin{workedexamplebox}{Reducing a Large Power}

Find the remainder when

\[ 7^4 \]

is divided by \(5\).

\begin{tabularx}{\textwidth}{@{}>{\raggedright\arraybackslash}p{0.42\textwidth}X@{}}
Reduce \(7\) modulo \(5\). &
\(7\equiv2\pmod5\)\\
Raise both sides to the fourth power. &
\(7^4\equiv2^4\pmod5\)\\
Compute the smaller power. &
\(2^4=16\)\\
Reduce again. &
\(16\equiv1\pmod5\)
\end{tabularx}

\begin{answerbox}
\[ \boxed{7^4\equiv1\pmod5.} \]
The remainder is \(1\).
\end{answerbox}
\end{workedexamplebox}

%%%%%%%%%%%%%%%%%%%%%%%%%%%%%%%%%%%%%%%%%%%%%%%%%%%%%%%%%%%%
\subsection{Modular Addition and Multiplication}
\label{subsec:modular-addition-multiplication}
%%%%%%%%%%%%%%%%%%%%%%%%%%%%%%%%%%%%%%%%%%%%%%%%%%%%%%%%%%%%

Because congruence is compatible with addition and multiplication, arithmetic
may be performed directly with residue classes.

For example, modulo \(7\),

\[ 5+6=11\equiv4\pmod7. \]

Similarly,

\[ 5(6)=30\equiv2\pmod7. \]

Thus in \(\Z/7\Z\),

\[ [5]+[6]=[4] \]

and

\[ [5][6]=[2]. \]

%%%%%%%%%%%%%%%%%%%%%%%%%%%%%%%%%%%%%%%%%%%%%%%%%%%%%%%%%%%%
\subsection{Cancellation in Congruences}
\label{subsec:congruence-cancellation}
%%%%%%%%%%%%%%%%%%%%%%%%%%%%%%%%%%%%%%%%%%%%%%%%%%%%%%%%%%%%

Ordinary cancellation must be used carefully in modular arithmetic.

From

\[ ac\equiv bc\pmod n, \]

we cannot always conclude

\[ a\equiv b\pmod n. \]

For example,

\[ 2(1)\equiv2(4)\pmod6 \]

because

\[ 2\equiv8\pmod6. \]

But

\[ 1\not\equiv4\pmod6. \]

Cancellation is valid when the cancelled factor is relatively prime to the
modulus.

\begin{theorem}[Modular Cancellation]
If

\[ ac\equiv bc\pmod n \]

and

\[ \gcd(c,n)=1, \]

then

\[ a\equiv b\pmod n. \]
\end{theorem}

%%%%%%%%%%%%%%%%%%%%%%%%%%%%%%%%%%%%%%%%%%%%%%%%%%%%%%%%%%%%
\subsection{Linear Congruences}
\label{subsec:linear-congruences}
%%%%%%%%%%%%%%%%%%%%%%%%%%%%%%%%%%%%%%%%%%%%%%%%%%%%%%%%%%%%

A \emph{linear congruence} has the form

\[ ax\equiv b\pmod n. \]

The question is whether there exists an integer \(x\) satisfying the
congruence.

\begin{theorem}[Solvability of a Linear Congruence]
The congruence

\[ ax\equiv b\pmod n \]

has a solution if and only if

\[ \gcd(a,n)\mid b. \]
\end{theorem}

If

\[ d=\gcd(a,n) \]

and \(d\mid b\), then there are exactly \(d\) incongruent solutions modulo
\(n\).

%%%%%%%%%%%%%%%%%%%%%%%%%%%%%%%%%%%%%%%%%%%%%%%%%%%%%%%%%%%%
\subsection{Worked Example: Linear Congruence}
\label{subsec:worked-linear-congruence}
%%%%%%%%%%%%%%%%%%%%%%%%%%%%%%%%%%%%%%%%%%%%%%%%%%%%%%%%%%%%

\begin{workedexamplebox}{Solving a Linear Congruence}

Solve

\[ 3x\equiv4\pmod7. \]

\begin{tabularx}{\textwidth}{@{}>{\raggedright\arraybackslash}p{0.42\textwidth}X@{}}
Find an inverse of \(3\) modulo \(7\). &
\(3(5)=15\equiv1\pmod7\)\\
Multiply both sides by \(5\). &
\(15x\equiv20\pmod7\)\\
Reduce modulo \(7\). &
\(x\equiv6\pmod7\)
\end{tabularx}

\begin{answerbox}
\[ \boxed{x\equiv6\pmod7.} \]
\end{answerbox}
\end{workedexamplebox}

%%%%%%%%%%%%%%%%%%%%%%%%%%%%%%%%%%%%%%%%%%%%%%%%%%%%%%%%%%%%
\subsection{Multiplicative Inverses Modulo \(n\)}
\label{subsec:modular-inverses}
%%%%%%%%%%%%%%%%%%%%%%%%%%%%%%%%%%%%%%%%%%%%%%%%%%%%%%%%%%%%

\begin{definitionbox}{Multiplicative Inverse Modulo \(n\)}
An integer \(b\) is a \emph{multiplicative inverse} of \(a\) modulo \(n\) if

\[ ab\equiv1\pmod n. \]
\end{definitionbox}

Such an inverse exists precisely when

\[ \boxed{\gcd(a,n)=1.} \]

This follows from Bézout's Identity.

If

\[ ax+ny=1, \]

then

\[ ax\equiv1\pmod n. \]

Thus \(x\) is an inverse of \(a\) modulo \(n\).

%%%%%%%%%%%%%%%%%%%%%%%%%%%%%%%%%%%%%%%%%%%%%%%%%%%%%%%%%%%%
\subsection{Worked Example: Modular Inverse}
\label{subsec:worked-modular-inverse}
%%%%%%%%%%%%%%%%%%%%%%%%%%%%%%%%%%%%%%%%%%%%%%%%%%%%%%%%%%%%

\begin{workedexamplebox}{Finding an Inverse Modulo \(26\)}

Find the multiplicative inverse of \(7\) modulo \(26\).

\begin{tabularx}{\textwidth}{@{}>{\raggedright\arraybackslash}p{0.42\textwidth}X@{}}
Apply the Euclidean Algorithm. &
\(26=7(3)+5\)\\
Continue. &
\(7=5(1)+2\)\\
Continue. &
\(5=2(2)+1\)\\
Work backward. &
\(1=5-2(2)\)\\
Replace \(2=7-5\). &
\(1=3(5)-2(7)\)\\
Replace \(5=26-3(7)\). &
\(1=3(26)-11(7)\)\\
Reduce modulo \(26\). &
\(-11(7)\equiv1\pmod{26}\)
\end{tabularx}

Since

\[ -11\equiv15\pmod{26}, \]

\begin{answerbox}
\[ \boxed{7^{-1}\equiv15\pmod{26}.} \]
\end{answerbox}
\end{workedexamplebox}

%%%%%%%%%%%%%%%%%%%%%%%%%%%%%%%%%%%%%%%%%%%%%%%%%%%%%%%%%%%%
\subsection{Systems of Congruences}
\label{subsec:systems-congruences}
%%%%%%%%%%%%%%%%%%%%%%%%%%%%%%%%%%%%%%%%%%%%%%%%%%%%%%%%%%%%

Number theory often requires an integer to satisfy several congruence
conditions simultaneously.

For example,

\[ x\equiv2\pmod3, \]

\[ x\equiv3\pmod5. \]

One solution is

\[ x=8. \]

Indeed,

\[ 8\equiv2\pmod3 \]

and

\[ 8\equiv3\pmod5. \]

All other solutions are congruent to \(8\) modulo \(15\).

%%%%%%%%%%%%%%%%%%%%%%%%%%%%%%%%%%%%%%%%%%%%%%%%%%%%%%%%%%%%
\subsection{The Chinese Remainder Theorem}
\label{subsec:chinese-remainder-theorem}
%%%%%%%%%%%%%%%%%%%%%%%%%%%%%%%%%%%%%%%%%%%%%%%%%%%%%%%%%%%%

\begin{theorem}[Chinese Remainder Theorem]
Suppose

\[ n_1,n_2,\ldots,n_k \]

are pairwise relatively prime positive integers.

Then the system

\[
x\equiv a_1\pmod{n_1},\qquad
x\equiv a_2\pmod{n_2},\qquad
\ldots,\qquad
x\equiv a_k\pmod{n_k}
\]

has a solution.

Moreover, the solution is unique modulo

\[ N=n_1n_2\cdots n_k. \]
\end{theorem}

The theorem allows several smaller modular conditions to be combined into a
single congruence.

%%%%%%%%%%%%%%%%%%%%%%%%%%%%%%%%%%%%%%%%%%%%%%%%%%%%%%%%%%%%
\subsection{Worked Example: Chinese Remainder Theorem}
\label{subsec:worked-crt}
%%%%%%%%%%%%%%%%%%%%%%%%%%%%%%%%%%%%%%%%%%%%%%%%%%%%%%%%%%%%

\begin{workedexamplebox}{Solving a Congruence System}

Solve

\[ x\equiv2\pmod3,\qquad x\equiv3\pmod5. \]

\begin{tabularx}{\textwidth}{@{}>{\raggedright\arraybackslash}p{0.42\textwidth}X@{}}
Use the first congruence. &
\(x=2+3k\)\\
Substitute into the second. &
\(2+3k\equiv3\pmod5\)\\
Simplify. &
\(3k\equiv1\pmod5\)\\
The inverse of \(3\) modulo \(5\) is \(2\). &
\(k\equiv2\pmod5\)\\
Take \(k=2\). &
\(x=2+3(2)=8\)
\end{tabularx}

Since the moduli are relatively prime, the solution is unique modulo

\[ 3(5)=15. \]

\begin{answerbox}
\[ \boxed{x\equiv8\pmod{15}.} \]
\end{answerbox}
\end{workedexamplebox}

%%%%%%%%%%%%%%%%%%%%%%%%%%%%%%%%%%%%%%%%%%%%%%%%%%%%%%%%%%%%
\subsection{Powers Modulo \(n\)}
\label{subsec:powers-modulo-n}
%%%%%%%%%%%%%%%%%%%%%%%%%%%%%%%%%%%%%%%%%%%%%%%%%%%%%%%%%%%%

Large powers often become manageable because residues repeat.

For example, consider powers of \(2\) modulo \(5\):

\[ 2^1\equiv2,\qquad2^2\equiv4,\qquad2^3\equiv3,\qquad2^4\equiv1\pmod5. \]

The cycle then repeats because

\[ 2^{k+4}\equiv2^k\pmod5. \]

Therefore

\[ 2^{100}\equiv1\pmod5. \]

Theorems of Fermat and Euler generalize this phenomenon.

%%%%%%%%%%%%%%%%%%%%%%%%%%%%%%%%%%%%%%%%%%%%%%%%%%%%%%%%%%%%
\subsection{Fermat's Little Theorem}
\label{subsec:fermats-little-theorem}
%%%%%%%%%%%%%%%%%%%%%%%%%%%%%%%%%%%%%%%%%%%%%%%%%%%%%%%%%%%%

\begin{theorem}[Fermat's Little Theorem]
If \(p\) is prime and

\[ p\nmid a, \]

then

\[ a^{p-1}\equiv1\pmod p. \]
\end{theorem}

An equivalent form is

\[ a^p\equiv a\pmod p \]

for every integer \(a\).

\begin{importantbox}
Fermat's Little Theorem is useful for reducing large powers modulo a prime
number.
\end{importantbox}

%%%%%%%%%%%%%%%%%%%%%%%%%%%%%%%%%%%%%%%%%%%%%%%%%%%%%%%%%%%%
\subsection{Worked Example: Fermat's Little Theorem}
\label{subsec:worked-fermat-little}
%%%%%%%%%%%%%%%%%%%%%%%%%%%%%%%%%%%%%%%%%%%%%%%%%%%%%%%%%%%%

\begin{workedexamplebox}{Reducing a Large Power}

Find

\[ 3^{100}\pmod7. \]

\begin{tabularx}{\textwidth}{@{}>{\raggedright\arraybackslash}p{0.42\textwidth}X@{}}
Apply Fermat's Little Theorem. &
\(3^6\equiv1\pmod7\)\\
Divide the exponent by \(6\). &
\(100=6(16)+4\)\\
Rewrite the power. &
\(3^{100}=(3^6)^{16}3^4\)\\
Reduce modulo \(7\). &
\(3^{100}\equiv3^4=81\equiv4\pmod7\)
\end{tabularx}

\begin{answerbox}
\[ \boxed{3^{100}\equiv4\pmod7.} \]
\end{answerbox}
\end{workedexamplebox}

%%%%%%%%%%%%%%%%%%%%%%%%%%%%%%%%%%%%%%%%%%%%%%%%%%%%%%%%%%%%
\subsection{Euler's Phi Function}
\label{subsec:euler-phi-function}
%%%%%%%%%%%%%%%%%%%%%%%%%%%%%%%%%%%%%%%%%%%%%%%%%%%%%%%%%%%%

\begin{definitionbox}{Euler's Phi Function}
For a positive integer \(n\), Euler's phi function

\[ \varphi(n) \]

is the number of integers in

\[ \{1,2,\ldots,n\} \]

that are relatively prime to \(n\).
\end{definitionbox}

The Greek letter

\[ \varphi \]

is a form of the Greek letter \emph{phi}, commonly pronounced ``fie'' or
``fee'' depending on convention.

The notation

\[ \varphi(n) \]

is read ``phi of \(n\).''

For example,

\[ \varphi(8)=4 \]

because

\[ 1,3,5,7 \]

are the integers from \(1\) through \(8\) relatively prime to \(8\).

%%%%%%%%%%%%%%%%%%%%%%%%%%%%%%%%%%%%%%%%%%%%%%%%%%%%%%%%%%%%
\subsection{Phi of a Prime}
\label{subsec:phi-prime}
%%%%%%%%%%%%%%%%%%%%%%%%%%%%%%%%%%%%%%%%%%%%%%%%%%%%%%%%%%%%

If \(p\) is prime, every integer

\[ 1,2,\ldots,p-1 \]

is relatively prime to \(p\).

Therefore

\[ \boxed{\varphi(p)=p-1.} \]

For a prime power,

\[ \boxed{\varphi(p^k)=p^k-p^{k-1}=p^k\left(1-\frac1p\right).} \]

%%%%%%%%%%%%%%%%%%%%%%%%%%%%%%%%%%%%%%%%%%%%%%%%%%%%%%%%%%%%
\subsection{Euler Phi Product Formula}
\label{subsec:phi-product-formula}
%%%%%%%%%%%%%%%%%%%%%%%%%%%%%%%%%%%%%%%%%%%%%%%%%%%%%%%%%%%%

If

\[ n=p_1^{\alpha_1}\cdots p_k^{\alpha_k} \]

is the prime factorization of \(n\), then

\[ \boxed{\varphi(n)=n\prod_{p\mid n}\left(1-\frac1p\right).} \]

For example,

\[ 12=2^2\cdot3, \]

so

\[ \varphi(12)=12\left(1-\frac12\right)\left(1-\frac13\right)=4. \]

Indeed, the integers

\[ 1,5,7,11 \]

are relatively prime to \(12\).

%%%%%%%%%%%%%%%%%%%%%%%%%%%%%%%%%%%%%%%%%%%%%%%%%%%%%%%%%%%%
\subsection{Euler's Theorem}
\label{subsec:eulers-theorem-number-theory}
%%%%%%%%%%%%%%%%%%%%%%%%%%%%%%%%%%%%%%%%%%%%%%%%%%%%%%%%%%%%

\begin{theorem}[Euler's Theorem]
If

\[ \gcd(a,n)=1, \]

then

\[ a^{\varphi(n)}\equiv1\pmod n. \]
\end{theorem}

Fermat's Little Theorem is the prime-modulus case because

\[ \varphi(p)=p-1. \]

%%%%%%%%%%%%%%%%%%%%%%%%%%%%%%%%%%%%%%%%%%%%%%%%%%%%%%%%%%%%
\subsection{Worked Example: Euler's Theorem}
\label{subsec:worked-eulers-theorem}
%%%%%%%%%%%%%%%%%%%%%%%%%%%%%%%%%%%%%%%%%%%%%%%%%%%%%%%%%%%%

\begin{workedexamplebox}{Large Powers Modulo a Composite Number}

Find

\[ 7^{100}\pmod{15}. \]

\begin{tabularx}{\textwidth}{@{}>{\raggedright\arraybackslash}p{0.42\textwidth}X@{}}
Check relative primality. &
\(\gcd(7,15)=1\)\\
Compute the phi value. &
\(\varphi(15)=15(1-\frac13)(1-\frac15)=8\)\\
Apply Euler's Theorem. &
\(7^8\equiv1\pmod{15}\)\\
Reduce the exponent. &
\(100=8(12)+4\)\\
Therefore &
\(7^{100}\equiv7^4\pmod{15}\)\\
Compute. &
\(7^2=49\equiv4\pmod{15}\)\\
Square again. &
\(7^4\equiv4^2=16\equiv1\pmod{15}\)
\end{tabularx}

\begin{answerbox}
\[ \boxed{7^{100}\equiv1\pmod{15}.} \]
\end{answerbox}
\end{workedexamplebox}

%%%%%%%%%%%%%%%%%%%%%%%%%%%%%%%%%%%%%%%%%%%%%%%%%%%%%%%%%%%%
\subsection{Diophantine Equations}
\label{subsec:diophantine-equations-elementary}
%%%%%%%%%%%%%%%%%%%%%%%%%%%%%%%%%%%%%%%%%%%%%%%%%%%%%%%%%%%%

A central theme of number theory is solving equations under the requirement
that the solutions be integers.

\begin{definitionbox}{Diophantine Equation}
A \emph{Diophantine equation} is an equation for which integer solutions are
sought.
\end{definitionbox}

The name comes from the ancient mathematician Diophantus.

Examples include

\[ 3x+5y=17, \]

\[ x^2+y^2=z^2, \]

and

\[ x^2-2y^2=1. \]

Unlike equations over \(\R\), a Diophantine equation may have no integer
solution even when real solutions exist.

%%%%%%%%%%%%%%%%%%%%%%%%%%%%%%%%%%%%%%%%%%%%%%%%%%%%%%%%%%%%
\subsection{Linear Diophantine Equations}
\label{subsec:linear-diophantine-equations}
%%%%%%%%%%%%%%%%%%%%%%%%%%%%%%%%%%%%%%%%%%%%%%%%%%%%%%%%%%%%

The simplest form is

\[ ax+by=c, \]

where

\[ a,b,c\in\Z. \]

\begin{theorem}[Linear Diophantine Solvability]
The equation

\[ ax+by=c \]

has an integer solution if and only if

\[ \gcd(a,b)\mid c. \]
\end{theorem}

This result follows directly from Bézout's Identity.

%%%%%%%%%%%%%%%%%%%%%%%%%%%%%%%%%%%%%%%%%%%%%%%%%%%%%%%%%%%%
\subsection{Worked Example: Linear Diophantine Equation}
\label{subsec:worked-linear-diophantine}
%%%%%%%%%%%%%%%%%%%%%%%%%%%%%%%%%%%%%%%%%%%%%%%%%%%%%%%%%%%%

\begin{workedexamplebox}{Finding Integer Solutions}

Solve

\[ 15x+21y=6. \]

\begin{tabularx}{\textwidth}{@{}>{\raggedright\arraybackslash}p{0.42\textwidth}X@{}}
Compute the gcd. &
\(\gcd(15,21)=3\)\\
Check solvability. &
\(3\mid6\), so solutions exist.\\
Divide the equation by \(3\). &
\(5x+7y=2\)\\
Find one convenient solution. &
\(x=-1,\ y=1\)\\
Verify. &
\(5(-1)+7(1)=2\)
\end{tabularx}

Thus one solution of the original equation is

\[ (x,y)=(-1,1). \]

\begin{answerbox}
\[ \boxed{x=-1,\qquad y=1} \]
is one integer solution.
\end{answerbox}
\end{workedexamplebox}

%%%%%%%%%%%%%%%%%%%%%%%%%%%%%%%%%%%%%%%%%%%%%%%%%%%%%%%%%%%%
\subsection{All Solutions of a Linear Diophantine Equation}
\label{subsec:all-linear-diophantine-solutions}
%%%%%%%%%%%%%%%%%%%%%%%%%%%%%%%%%%%%%%%%%%%%%%%%%%%%%%%%%%%%

Suppose

\[ ax+by=c \]

has one solution

\[ (x_0,y_0). \]

Let

\[ d=\gcd(a,b). \]

Then every integer solution has the form

\[ \boxed{x=x_0+\frac{b}{d}t,\qquad y=y_0-\frac{a}{d}t,\qquad t\in\Z.} \]

For

\[ 15x+21y=6, \]

we have

\[ d=3 \]

and one solution is

\[ (x_0,y_0)=(-1,1). \]

Therefore all solutions are

\[ \boxed{x=-1+7t,\qquad y=1-5t,\qquad t\in\Z.} \]

%%%%%%%%%%%%%%%%%%%%%%%%%%%%%%%%%%%%%%%%%%%%%%%%%%%%%%%%%%%%
\subsection{Pythagorean Triples}
\label{subsec:pythagorean-triples}
%%%%%%%%%%%%%%%%%%%%%%%%%%%%%%%%%%%%%%%%%%%%%%%%%%%%%%%%%%%%

One of the most famous Diophantine equations is

\[ x^2+y^2=z^2. \]

An integer solution with positive entries is called a
\emph{Pythagorean triple}.

Examples include

\[ (3,4,5),\qquad(5,12,13),\qquad(8,15,17). \]

A primitive Pythagorean triple satisfies

\[ \gcd(x,y,z)=1. \]

Every primitive Pythagorean triple can be generated, up to exchanging the
legs, using

\[ x=m^2-n^2,\qquad y=2mn,\qquad z=m^2+n^2, \]

where

\[ m>n>0,\qquad \gcd(m,n)=1, \]

and \(m,n\) have opposite parity.

%%%%%%%%%%%%%%%%%%%%%%%%%%%%%%%%%%%%%%%%%%%%%%%%%%%%%%%%%%%%
\subsection{Divisor Functions}
\label{subsec:divisor-functions}
%%%%%%%%%%%%%%%%%%%%%%%%%%%%%%%%%%%%%%%%%%%%%%%%%%%%%%%%%%%%

Prime factorization also allows arithmetic information about divisors to be
computed efficiently.

Suppose

\[ n=p_1^{\alpha_1}p_2^{\alpha_2}\cdots p_k^{\alpha_k}. \]

Every positive divisor of \(n\) has the form

\[ p_1^{\beta_1}p_2^{\beta_2}\cdots p_k^{\beta_k}, \]

where

\[ 0\leq\beta_i\leq\alpha_i. \]

Therefore the number of positive divisors is

\[ \boxed{\tau(n)=\prod_{i=1}^{k}(\alpha_i+1).} \]

The Greek letter

\[ \tau \]

is called \emph{tau}, pronounced ``tow.''

Here \(\tau(n)\) denotes the number-of-divisors function.

%%%%%%%%%%%%%%%%%%%%%%%%%%%%%%%%%%%%%%%%%%%%%%%%%%%%%%%%%%%%
\subsection{Worked Example: Counting Divisors}
\label{subsec:worked-counting-divisors}
%%%%%%%%%%%%%%%%%%%%%%%%%%%%%%%%%%%%%%%%%%%%%%%%%%%%%%%%%%%%

\begin{workedexamplebox}{Number of Positive Divisors}

How many positive divisors does \(360\) have?

\begin{tabularx}{\textwidth}{@{}>{\raggedright\arraybackslash}p{0.42\textwidth}X@{}}
Prime factorize \(360\). &
\(360=2^3\cdot3^2\cdot5^1\)\\
Add \(1\) to each exponent. &
\(3+1=4,\quad2+1=3,\quad1+1=2\)\\
Multiply. &
\(\tau(360)=4(3)(2)=24\)
\end{tabularx}

\begin{answerbox}
\[ \boxed{\tau(360)=24.} \]
\end{answerbox}
\end{workedexamplebox}

%%%%%%%%%%%%%%%%%%%%%%%%%%%%%%%%%%%%%%%%%%%%%%%%%%%%%%%%%%%%
\subsection{Sum of Divisors}
\label{subsec:sum-divisors}
%%%%%%%%%%%%%%%%%%%%%%%%%%%%%%%%%%%%%%%%%%%%%%%%%%%%%%%%%%%%

Another arithmetic function records the sum of positive divisors.

It is commonly written

\[ \sigma(n). \]

The Greek letter

\[ \sigma \]

is called \emph{sigma}.

For

\[ n=p_1^{\alpha_1}\cdots p_k^{\alpha_k}, \]

the divisor-sum formula is

\[
\boxed{
\sigma(n)
=
\prod_{i=1}^{k}
\left(
1+p_i+p_i^2+\cdots+p_i^{\alpha_i}
\right).
}
\]

Using the finite geometric-series formula,

\[
\boxed{
\sigma(n)
=
\prod_{i=1}^{k}
\frac{p_i^{\alpha_i+1}-1}{p_i-1}.
}
\]

%%%%%%%%%%%%%%%%%%%%%%%%%%%%%%%%%%%%%%%%%%%%%%%%%%%%%%%%%%%%
\subsection{Perfect Numbers}
\label{subsec:perfect-numbers}
%%%%%%%%%%%%%%%%%%%%%%%%%%%%%%%%%%%%%%%%%%%%%%%%%%%%%%%%%%%%

A positive integer is called \emph{perfect} when it equals the sum of its
positive proper divisors.

For example, the proper divisors of \(6\) are

\[ 1,2,3. \]

Since

\[ 1+2+3=6, \]

the number \(6\) is perfect.

Equivalently,

\[ \sigma(n)=2n. \]

Other examples include

\[ 28,\qquad496,\qquad8128. \]

Perfect numbers provide an early example of how elementary questions about
divisors lead to deeper number-theoretic problems.

%%%%%%%%%%%%%%%%%%%%%%%%%%%%%%%%%%%%%%%%%%%%%%%%%%%%%%%%%%%%
\subsection{Quadratic Residues}
\label{subsec:quadratic-residues-introduction}
%%%%%%%%%%%%%%%%%%%%%%%%%%%%%%%%%%%%%%%%%%%%%%%%%%%%%%%%%%%%

Modular arithmetic can also be used to study squares.

\begin{definitionbox}{Quadratic Residue}
Let \(n>1\). An integer \(a\) is called a \emph{quadratic residue modulo}
\(n\) if the congruence

\[ x^2\equiv a\pmod n \]

has an integer solution.
\end{definitionbox}

For example, modulo \(7\),

\[ 1^2\equiv1,\qquad2^2\equiv4,\qquad3^2\equiv2\pmod7. \]

Thus the nonzero quadratic residues modulo \(7\) are

\[ 1,2,4. \]

%%%%%%%%%%%%%%%%%%%%%%%%%%%%%%%%%%%%%%%%%%%%%%%%%%%%%%%%%%%%
\subsection{Quadratic Residue Table}
\label{subsec:quadratic-residue-table}
%%%%%%%%%%%%%%%%%%%%%%%%%%%%%%%%%%%%%%%%%%%%%%%%%%%%%%%%%%%%

Squaring each residue modulo \(7\) gives:

\begin{center}
\begin{tabular}{c|ccccccc}
\toprule
\(x\) & \(0\) & \(1\) & \(2\) & \(3\) & \(4\) & \(5\) & \(6\)\\
\midrule
\(x^2\bmod7\) & \(0\) & \(1\) & \(4\) & \(2\) & \(2\) & \(4\) & \(1\)\\
\bottomrule
\end{tabular}
\end{center}

Therefore the possible square residues are

\[ 0,1,2,4. \]

Quadratic residues will become much more important in later number theory,
particularly in quadratic reciprocity and algebraic number theory.

%%%%%%%%%%%%%%%%%%%%%%%%%%%%%%%%%%%%%%%%%%%%%%%%%%%%%%%%%%%%
\subsection{Number Theory as an Algorithmic Subject}
\label{subsec:number-theory-algorithms}
%%%%%%%%%%%%%%%%%%%%%%%%%%%%%%%%%%%%%%%%%%%%%%%%%%%%%%%%%%%%

Elementary number theory contains many naturally computational procedures.

Examples include:

\begin{itemize}
    \item the Division Algorithm;
    \item the Euclidean Algorithm;
    \item the Extended Euclidean Algorithm;
    \item prime factorization;
    \item modular exponentiation;
    \item primality testing;
    \item solving linear congruences;
    \item solving systems of congruences.
\end{itemize}

\begin{centereddiagram}
\begin{tikzpicture}[
    node distance=9mm,
    box/.style={draw,rounded corners,align=center,minimum width=4.5cm,
    minimum height=0.9cm}
]
\node[box] (div) {Divisibility};
\node[box,below=of div] (gcd) {Greatest Common Divisors};
\node[box,below=of gcd] (euclid) {Euclidean Algorithm};
\node[box,below=of euclid] (bezout) {Bézout's Identity};
\node[box,below=of bezout] (inverse) {Modular Inverses};
\node[box,below=of inverse] (cong) {Congruences and Cryptography};

\draw[-{Latex[length=2.5mm]}] (div)--(gcd);
\draw[-{Latex[length=2.5mm]}] (gcd)--(euclid);
\draw[-{Latex[length=2.5mm]}] (euclid)--(bezout);
\draw[-{Latex[length=2.5mm]}] (bezout)--(inverse);
\draw[-{Latex[length=2.5mm]}] (inverse)--(cong);
\end{tikzpicture}
\end{centereddiagram}

This algorithmic character makes number theory particularly important in
computer science.

%%%%%%%%%%%%%%%%%%%%%%%%%%%%%%%%%%%%%%%%%%%%%%%%%%%%%%%%%%%%
\subsection{Elementary Cryptographic Connection}
\label{subsec:elementary-cryptographic-connection}
%%%%%%%%%%%%%%%%%%%%%%%%%%%%%%%%%%%%%%%%%%%%%%%%%%%%%%%%%%%%

Modern public-key cryptography uses arithmetic involving very large
integers.

Several elementary number-theoretic ideas appear directly:

\begin{itemize}
    \item prime numbers;
    \item greatest common divisors;
    \item modular inverses;
    \item Euler's phi function;
    \item modular exponentiation;
    \item Fermat's and Euler's theorems.
\end{itemize}

For example, RSA cryptography relies on the practical difficulty of
factoring a sufficiently large integer formed from large prime factors.

The cryptographic theory itself belongs to more advanced applications, but
its arithmetic foundation begins here.

%%%%%%%%%%%%%%%%%%%%%%%%%%%%%%%%%%%%%%%%%%%%%%%%%%%%%%%%%%%%
\subsection{Common Errors}
\label{subsec:elementary-number-theory-common-errors}
%%%%%%%%%%%%%%%%%%%%%%%%%%%%%%%%%%%%%%%%%%%%%%%%%%%%%%%%%%%%

\subsubsection{Reversing Divisibility}

The statement

\[ a\mid b \]

means \(b\) is a multiple of \(a\).

It does not mean that \(a\) is a multiple of \(b\).

%%%%%%%%%%%%%%%%%%%%%%%%%%%%%%%%%%%%%%%%%%%%%%%%%%%%%%%%%%%%
\subsubsection{Treating Divisibility Like a Fraction}

The notation

\[ a\mid b \]

is a mathematical relation.

It is not the fraction

\[ \frac ab. \]

%%%%%%%%%%%%%%%%%%%%%%%%%%%%%%%%%%%%%%%%%%%%%%%%%%%%%%%%%%%%
\subsubsection{Allowing an Invalid Remainder}

In

\[ a=dq+r, \]

the Division Algorithm requires

\[ 0\leq r<d. \]

A negative remainder does not satisfy the standard theorem.

%%%%%%%%%%%%%%%%%%%%%%%%%%%%%%%%%%%%%%%%%%%%%%%%%%%%%%%%%%%%
\subsubsection{Confusing GCD and LCM}

The gcd is the greatest positive integer dividing both numbers.

The lcm is the smallest positive integer divisible by both numbers.

%%%%%%%%%%%%%%%%%%%%%%%%%%%%%%%%%%%%%%%%%%%%%%%%%%%%%%%%%%%%
\subsubsection{Calling \(1\) Prime}

The number \(1\) is neither prime nor composite.

Prime numbers begin with

\[ 2,3,5,7,\ldots \]

%%%%%%%%%%%%%%%%%%%%%%%%%%%%%%%%%%%%%%%%%%%%%%%%%%%%%%%%%%%%
\subsubsection{Cancelling Illegally Modulo \(n\)}

From

\[ ac\equiv bc\pmod n, \]

one may cancel \(c\) directly only when \(c\) is invertible modulo \(n\),
equivalently when

\[ \gcd(c,n)=1. \]

%%%%%%%%%%%%%%%%%%%%%%%%%%%%%%%%%%%%%%%%%%%%%%%%%%%%%%%%%%%%
\subsubsection{Assuming Every Nonzero Residue Has an Inverse}

An integer \(a\) has an inverse modulo \(n\) precisely when

\[ \gcd(a,n)=1. \]

For example, \(2\) has no multiplicative inverse modulo \(6\).

%%%%%%%%%%%%%%%%%%%%%%%%%%%%%%%%%%%%%%%%%%%%%%%%%%%%%%%%%%%%
\subsubsection{Using Fermat's Little Theorem with a Composite Modulus}

Fermat's Little Theorem specifically uses a prime modulus.

Euler's Theorem provides a more general result when

\[ \gcd(a,n)=1. \]

%%%%%%%%%%%%%%%%%%%%%%%%%%%%%%%%%%%%%%%%%%%%%%%%%%%%%%%%%%%%
\subsubsection{Forgetting the GCD Condition in a Diophantine Equation}

The equation

\[ ax+by=c \]

has integer solutions only when

\[ \gcd(a,b)\mid c. \]

%%%%%%%%%%%%%%%%%%%%%%%%%%%%%%%%%%%%%%%%%%%%%%%%%%%%%%%%%%%%
\subsubsection{Assuming Terms That Look Simple Have Simple Behavior}

Elementary-looking integer problems can be extremely difficult.

The simplicity of the definitions in number theory should not be confused
with simplicity of the subject.

%%%%%%%%%%%%%%%%%%%%%%%%%%%%%%%%%%%%%%%%%%%%%%%%%%%%%%%%%%%%
\subsection{Applications and Connections}
\label{subsec:elementary-number-theory-applications}
%%%%%%%%%%%%%%%%%%%%%%%%%%%%%%%%%%%%%%%%%%%%%%%%%%%%%%%%%%%%

\begin{applicationbox}
\textbf{Algebra.}
Congruence classes lead naturally to quotient rings such as
\(\Z/n\Z\), while modular inverses correspond to units in these rings.

\medskip

\textbf{Cryptography.}
Prime numbers, modular exponentiation, Euler's theorem, and modular inverses
form the arithmetic foundation of several cryptographic systems.

\medskip

\textbf{Computer Science.}
The Euclidean Algorithm is an important example of an efficient algorithm.
Modular arithmetic appears in hashing, checksums, random-number generation,
and computational algebra.

\medskip

\textbf{Coding Theory.}
Arithmetic over finite modular structures contributes to the construction
and analysis of error-detecting and error-correcting codes.

\medskip

\textbf{Geometry.}
Pythagorean triples connect integer arithmetic with right triangles and
Diophantine geometry.

\medskip

\textbf{Combinatorics.}
Divisibility, modular arguments, and arithmetic functions appear throughout
enumerative and algebraic combinatorics.

\medskip

\textbf{Probability.}
Arithmetic properties of integers motivate probabilistic questions about
prime distribution and random integers.

\medskip

\textbf{Advanced Number Theory.}
Prime distribution leads toward Analytic Number Theory, while extensions of
factorization and congruence lead toward Algebraic Number Theory.
\end{applicationbox}

%%%%%%%%%%%%%%%%%%%%%%%%%%%%%%%%%%%%%%%%%%%%%%%%%%%%%%%%%%%%
\subsection{Exercises}
\label{subsec:elementary-number-theory-exercises}
%%%%%%%%%%%%%%%%%%%%%%%%%%%%%%%%%%%%%%%%%%%%%%%%%%%%%%%%%%%%

\begin{exercise}[1]
Determine whether each statement is true:

\[ 4\mid28,\qquad6\mid42,\qquad8\mid50,\qquad11\mid121. \]
\end{exercise}

\begin{exercise}[1]
List all positive divisors of \(48\).
\end{exercise}

\begin{exercise}[2]
Prove that if \(a\mid b\) and \(b\mid c\), then \(a\mid c\).
\end{exercise}

\begin{exercise}[2]
Prove that if \(a\mid b\) and \(a\mid c\), then

\[ a\mid(5b-3c). \]
\end{exercise}

\begin{exercise}[1]
Use the Division Algorithm to write \(157\) in the form

\[ 157=13q+r. \]
\end{exercise}

\begin{exercise}[2]
Write \(-43\) in the form

\[ -43=7q+r,\qquad0\leq r<7. \]
\end{exercise}

\begin{exercise}[1]
Find

\[ \gcd(84,126). \]
\end{exercise}

\begin{exercise}[1]
Find

\[ \operatorname{lcm}(84,126). \]
\end{exercise}

\begin{exercise}[2]
Use the Euclidean Algorithm to compute

\[ \gcd(414,662). \]
\end{exercise}

\begin{exercise}[3]
Use the Extended Euclidean Algorithm to express

\[ \gcd(414,662) \]

as an integer linear combination of \(414\) and \(662\).
\end{exercise}

\begin{exercise}[2]
Determine whether \(35\) and \(64\) are relatively prime.
\end{exercise}

\begin{exercise}[2]
Prove that if

\[ \gcd(a,b)=1 \]

and

\[ a\mid c,\qquad b\mid c, \]

then

\[ ab\mid c. \]
\end{exercise}

\begin{exercise}[1]
Determine which of the following are prime:

\[ 29,\qquad51,\qquad71,\qquad91,\qquad97. \]
\end{exercise}

\begin{exercise}[1]
Find the canonical prime factorization of

\[ 1260. \]
\end{exercise}

\begin{exercise}[2]
Use prime factorizations to compute the gcd and lcm of

\[ 540\qquad\text{and}\qquad756. \]
\end{exercise}

\begin{exercise}[2]
Explain why it is enough to test prime divisors up to \(\sqrt n\) when
testing \(n\) for primality.
\end{exercise}

\begin{exercise}[2]
Reproduce Euclid's proof that there are infinitely many primes.
\end{exercise}

\begin{exercise}[1]
Determine whether

\[ 47\equiv5\pmod7. \]
\end{exercise}

\begin{exercise}[1]
Find the least nonnegative residue of

\[ 137\pmod{12}. \]
\end{exercise}

\begin{exercise}[1]
Find the least nonnegative residue of

\[ -17\pmod5. \]
\end{exercise}

\begin{exercise}[2]
Compute

\[ 37^2\pmod{10}. \]
\end{exercise}

\begin{exercise}[2]
Compute

\[ 2^{20}\pmod7. \]
\end{exercise}

\begin{exercise}[2]
Solve

\[ 5x\equiv3\pmod7. \]
\end{exercise}

\begin{exercise}[2]
Determine whether

\[ 6x\equiv5\pmod9 \]

has a solution.
\end{exercise}

\begin{exercise}[2]
Solve

\[ 6x\equiv3\pmod9. \]
\end{exercise}

\begin{exercise}[2]
Find the multiplicative inverse of \(11\) modulo \(26\), if it exists.
\end{exercise}

\begin{exercise}[2]
Determine whether \(12\) has a multiplicative inverse modulo \(35\).
\end{exercise}

\begin{exercise}[2]
Solve the system

\[ x\equiv1\pmod3,\qquad x\equiv4\pmod5. \]
\end{exercise}

\begin{exercise}[3]
Solve

\[ x\equiv2\pmod3,\qquad x\equiv3\pmod5,\qquad x\equiv2\pmod7. \]
\end{exercise}

\begin{exercise}[2]
Use Fermat's Little Theorem to compute

\[ 5^{100}\pmod{13}. \]
\end{exercise}

\begin{exercise}[1]
Compute

\[ \varphi(7),\qquad\varphi(10),\qquad\varphi(12). \]
\end{exercise}

\begin{exercise}[2]
Compute

\[ \varphi(72). \]
\end{exercise}

\begin{exercise}[2]
Use Euler's Theorem to compute

\[ 11^{100}\pmod{20}. \]
\end{exercise}

\begin{exercise}[2]
Determine whether the Diophantine equation

\[ 12x+18y=7 \]

has integer solutions.
\end{exercise}

\begin{exercise}[2]
Find one integer solution of

\[ 14x+21y=35. \]
\end{exercise}

\begin{exercise}[3]
Find all integer solutions of

\[ 14x+21y=35. \]
\end{exercise}

\begin{exercise}[2]
Verify that

\[ (20,21,29) \]

is a Pythagorean triple.
\end{exercise}

\begin{exercise}[2]
Use the Pythagorean-triple formulas with

\[ m=4,\qquad n=1 \]

to generate a triple.
\end{exercise}

\begin{exercise}[2]
Compute the number of positive divisors of

\[ 720. \]
\end{exercise}

\begin{exercise}[2]
Compute

\[ \sigma(28). \]

Use the result to verify that \(28\) is perfect.
\end{exercise}

\begin{exercise}[2]
Find all quadratic residues modulo \(5\).
\end{exercise}

\begin{exercise}[2]
Find all quadratic residues modulo \(11\).
\end{exercise}

\begin{exercise}[3]
Determine whether

\[ x^2\equiv3\pmod7 \]

has a solution.
\end{exercise}

\begin{exercise}[3]
Prove that the square of every odd integer is congruent to \(1\) modulo \(8\).
\end{exercise}

\begin{exercise}[3]
Prove that if \(n\) is odd, then

\[ n^2\equiv1\pmod8 \]

need not hold for every odd-looking algebraic expression unless the
expression itself is known to be an odd integer. Explain the role of the
hypothesis carefully.
\end{exercise}

\begin{exercise}[3]
Show that if

\[ a\equiv b\pmod n, \]

then

\[ a^k\equiv b^k\pmod n \]

for every positive integer \(k\).
\end{exercise}

\begin{exercise}[3]
Explain the chain of ideas

\[
\text{Divisibility}
\longrightarrow
\text{GCD}
\longrightarrow
\text{Euclidean Algorithm}
\longrightarrow
\text{Bézout Identity}
\longrightarrow
\text{Modular Inverses}
\longrightarrow
\text{Congruences}.
\]
\end{exercise}

%%%%%%%%%%%%%%%%%%%%%%%%%%%%%%%%%%%%%%%%%%%%%%%%%%%%%%%%%%%%
\subsection{Conceptual Summary}
\label{subsec:elementary-number-theory-summary}
%%%%%%%%%%%%%%%%%%%%%%%%%%%%%%%%%%%%%%%%%%%%%%%%%%%%%%%%%%%%

Elementary number theory begins with divisibility.

The statement

\[ a\mid b \]

means that

\[ b=ak \]

for some integer \(k\).

The Division Algorithm expresses every integer relative to a positive divisor
in the form

\[ a=dq+r,\qquad0\leq r<d. \]

Greatest common divisors measure the common divisibility structure of two
integers.

The Euclidean Algorithm efficiently computes gcds, while the Extended
Euclidean Algorithm produces Bézout coefficients satisfying

\[ ax+by=\gcd(a,b). \]

Prime numbers are the multiplicative building blocks of positive integers.
The Fundamental Theorem of Arithmetic guarantees that every integer greater
than \(1\) has a unique prime factorization.

Congruence

\[ a\equiv b\pmod n \]

means that \(a\) and \(b\) differ by a multiple of \(n\).

Congruence arithmetic leads naturally to residue classes, modular inverses,
linear congruences, and systems of congruences.

The Chinese Remainder Theorem combines compatible modular conditions.

Fermat's Little Theorem and Euler's Theorem simplify large powers in modular
arithmetic.

Euler's phi function counts residue classes relatively prime to a modulus.

Diophantine equations ask for integer solutions to algebraic equations, with
linear Diophantine equations controlled by greatest common divisors.

Arithmetic functions such as

\[ \varphi(n),\qquad\tau(n),\qquad\sigma(n) \]

encode structural information about integers.

Together these ideas establish the elementary arithmetic foundation for the
more advanced branches of number theory.

%%%%%%%%%%%%%%%%%%%%%%%%%%%%%%%%%%%%%%%%%%%%%%%%%%%%%%%%%%%%
\subsection{Section Connections}
\label{subsec:elementary-number-theory-connections}
%%%%%%%%%%%%%%%%%%%%%%%%%%%%%%%%%%%%%%%%%%%%%%%%%%%%%%%%%%%%

\chapterconnection{
Elementary Number Theory begins the systematic study of integer arithmetic.
Its theory of divisibility, greatest common divisors, prime numbers, and
congruences connects directly with the algebraic structures introduced in
Chapter 1. Residue classes modulo \(n\) form quotient rings, modular inverses
correspond to units, and prime moduli lead naturally to finite fields.
Diophantine equations provide a bridge between arithmetic and geometry.
Prime numbers and arithmetic functions prepare the reader for Analytic
Number Theory, where analytic methods are used to study the distribution of
primes. Unique factorization, congruences, and extensions of the integers
lead toward Algebraic Number Theory. Algorithms such as the Euclidean
Algorithm and modular exponentiation lead toward Computational Number Theory.
Quadratic residues prepare the way for deeper reciprocity laws, while
Diophantine equations eventually connect number theory with algebraic
geometry and arithmetic geometry.
}

%%%%%%%%%%%%%%%%%%%%%%%%%%%%%%%%%%%%%%%%%%%%%%%%%%%%%%%%%%%%
% SECTION SOURCE TRACKING
%%%%%%%%%%%%%%%%%%%%%%%%%%%%%%%%%%%%%%%%%%%%%%%%%%%%%%%%%%%%

%------------------------------------------------------------
% 2.1 Elementary Number Theory
%------------------------------------------------------------
%
% Source basis:
%   - Standard undergraduate elementary number theory.
%
% Used for:
%   - divisibility;
%   - Division Algorithm;
%   - greatest common divisors and least common multiples;
%   - Euclidean and Extended Euclidean Algorithms;
%   - Bezout's Identity;
%   - relatively prime integers;
%   - Euclid's Lemma;
%   - prime and composite numbers;
%   - Fundamental Theorem of Arithmetic;
%   - elementary primality testing;
%   - congruences and residue classes;
%   - modular arithmetic and cancellation;
%   - linear congruences;
%   - modular inverses;
%   - Chinese Remainder Theorem;
%   - Fermat's Little Theorem;
%   - Euler's phi function and Euler's Theorem;
%   - linear Diophantine equations;
%   - Pythagorean triples;
%   - divisor-counting and divisor-sum functions;
%   - perfect numbers;
%   - introductory quadratic residues;
%   - computational and cryptographic connections.
%
% Notes:
%   - Basic integer arithmetic and number systems are taught
%     canonically in Chapter 0.
%   - Quotient rings, groups, fields, and algebraic structures
%     are taught canonically in Chapter 1.
%   - Prime-distribution theory is deferred to Analytic Number Theory.
%   - Algebraic integers, number fields, ideals, and advanced
%     factorization theory are deferred to Algebraic Number Theory.
%   - Advanced algorithms and computational complexity are deferred
%     to Computational Number Theory.
%   - Advanced quadratic residue theory and reciprocity are developed
%     in later number-theoretic sections.
%   - Diophantine geometry is deferred to its dedicated section.
%
%%%%%%%%%%%%%%%%%%%%%%%%%%%%%%%%%%%%%%%%%%%%%%%%%%%%%%%%%%%%